\subsection{A bounded-flip-energy Hamiltonian failing at every positive inverse temperature}
\label{threshold:fixed-soft-code}

Fix $L>0$ once and for all.  Let $\mathcal C_n=\{(-1)^bc_a:a,b\}$,
so $K:=|\mathcal C_n|=2n$, and define the temperature-independent
Hamiltonian
\[
 S_L(x)=\sum_{v\in\mathcal C_n}e^{Lx\cdot v},
 \qquad H_{n,L}(x)=-\log S_L(x),
 \qquad \mu_{n,\beta,L}(x)\propto S_L(x)^\beta.
 \tag{SC}
\]
An additive normalization such as $+\log(2n)$ in $H$ has no effect on
the Gibbs law.  The energy is a global many-body interaction.  The
following theorem makes no bounded-arity or bounded-degree claim.

\begin{theorem}[A fixed soft-code Hamiltonian]
\label{threshold:soft-code-theorem}
For each fixed $\beta,L>0$, put $\rho_\beta=\tanh(\beta L)$.  There
exist positive constants $C,c$, depending on $\beta,L$ but not $n$, such
that
\[
 \|\mu_{n,\beta,L}-\nu_{n,\rho_\beta}\|_{\rm TV}\le Ce^{-cn}.
 \tag{SC-TV}
\]
Consequently $\mathcal A(\mu_{\beta,L})$ fails for every fixed
$\beta>0$.  At $\beta=0$ it holds by uniform-input AA.  Moreover, for
every $n$ and $\beta\ge0$:
\begin{enumerate}
\item $|H(x)-H(x^{(i)})|\le2L$ and
      $\max H-\min H\le\log\cosh(Ln)\le Ln$.
\item All marginals on at most three distinct spins are exactly uniform;
      in particular $\Cov(X)=I$.
\item With $\varepsilon_\beta=(1+e^{2\beta L})^{-1}$,
      $\varepsilon_\beta\le\Pr(X_i=1\mid X_{-i})\le1-\varepsilon_\beta$
      for every conditioning.
\item In natural logarithms,
      $H_{\rm Shannon}(X)\ge n h(\varepsilon_\beta)$ and
      $\max_x\mu(x)\le(1-\varepsilon_\beta)^n$.
\end{enumerate}
Thus arbitrarily small single-spin energy changes, extensive total
energy, almost maximal entropy density, and these exact low-order
marginals do not guarantee full AA at any fixed positive inverse
temperature when global interaction order is unrestricted.
\end{theorem}

\begin{proof}[Proof of the total-variation comparison]
Write $\rho=\tanh(\beta L)$ and
\[
 A_v(x)=e^{Lx\cdot v},\qquad
 Q(x)=\sum_v A_v(x)^\beta,\qquad
 R(x)=\frac{S_L(x)^\beta}{Q(x)}.
\]
The law $q(x):=Q(x)/\sum_yQ(y)$ is exactly $\nu_{n,\rho}$:
each $v$ has the same partition function $(2\cosh(\beta L))^n$,
and conditioned on $v$ the spins are $X_z=v_z\eta_z$ with mean-$\rho$
independent signs.  The power-mean inequalities give the global bounds
\[
 \begin{array}{ll}
 1\le R(x)\le K^{\beta-1},&\beta\ge1,\\
 K^{\beta-1}\le R(x)\le1,&0<\beta\le1.
 \end{array}
 \tag{SC-R}
\]
In particular $|R-1|\le K^{(\beta-1)_+}$ everywhere, where
$u_+=\max(u,0)$.

Condition on the sampled mode $v$.  Every other signed character $w$
differs from $v$ on at least $n/2$ coordinates.  If $D=\{z:v_z\ne w_z\}$,
then
\[
 X\cdot v-X\cdot w=2\sum_{z\in D}\eta_z.
\]
Hoeffding's inequality, using $|D|\ge n/2$, gives
\[
 \Pr\{X\cdot v-X\cdot w<\rho n/2\mid v\}
 \le e^{-\rho^2n/16}.
\]
Let $G$ be the event that some mode exceeds every other score by at
least $d:=\rho n/2$.  A union bound therefore yields
\[
 q(G^c)\le(K-1)e^{-\rho^2n/16}.
 \tag{SC-bad}
\]

On $G$, let $M=\max_v A_v(x)$.  Then
\[
 S_L=M(1+u),\quad Q=M^\beta(1+v'),\quad
 0\le u\le u_n:=(K-1)e^{-L\rho n/2},\quad
 0\le v'\le v_n:=(K-1)e^{-\beta L\rho n/2}.
\]
For all sufficiently large $n$, $u_n\le1$.  The mean value theorem
gives
\[
 |R-1|\le\beta2^{(\beta-1)_+}u_n+v_n
 \quad\hbox{on }G.
\]
Consequently
\[
 \mathbb E_q|R-1|\le e_n:={}
 \beta2^{(\beta-1)_+}(K-1)e^{-L\rho n/2}
 +(K-1)e^{-\beta L\rho n/2}
 +K^{(\beta-1)_+}(K-1)e^{-\rho^2n/16}.
 \tag{SC-error}
\]
For fixed positive $\beta,L$ this is exponentially small, including
when $\beta<1$.  Finally,
$d\mu/dq=R/\mathbb E_qR$, and whenever $e_n<1$,
\[
 \|\mu-q\|_{\rm TV}
 =\frac{\mathbb E_q|R-\mathbb E_qR|}{2\mathbb E_qR}
 \le\frac{e_n}{1-e_n}.
\]
This proves (SC-TV); enlarging $C$ covers the finitely many smaller
$n$.  Theorem~\ref{threshold:noisy-query-theorem} and
Lemma~\ref{threshold:hard-component} then prove the query separation.
\end{proof}

\begin{proof}[Proof of the structural assertions]
A single flip multiplies each summand in $S_L$ by a factor in
$[e^{-2L},e^{2L}]$, proving the energy increment bound.  Pairing $v,-v$
shows
$2n\le S_L(x)\le2n\cosh(Ln)$, proving the energy-range bound.

Multiplication of $x$ coordinatewise by any signed character permutes
the modes in $S_L$, hence preserves $H$ and every Gibbs power.  For a
nonempty set $T$ of at most three distinct coordinates, one can choose
a signed character whose product on $T$ is $-1$: use the global sign
if $|T|$ is odd, and for two distinct $z,z'$ choose a character with
$a\cdot(z+z')=1$.  Thus
$\mathbb E\prod_{z\in T}X_z=0$.  Fourier inversion on each marginal
proves exact three-wise independence.

The energy increment bound makes every single-site odds ratio lie in
$[e^{-2\beta L},e^{2\beta L}]$, proving assertion (3).  The same
probability bounds hold when only an arbitrary subset of the other
spins is conditioned, by conditional expectation.  Applying the
chain rule for entropy gives the lower bound in (4); multiplying
successive conditional probability bounds gives the maximum-atom
bound.  For fixed $\beta$, choosing $L$ as small as desired makes
$h(\varepsilon_\beta)$ approach $\log2$ while the query obstruction
persists for every fixed $L>0$.
\end{proof}

At $\beta=1$, the comparison in (SC-TV) is an equality of laws, since
$S_L^1=Q$.  For other $\beta$, the theorem compares the actual Gibbs
power of the \emph{same} Hamiltonian with a noise model; it does not
silently replace $L$ by $\beta L$ in the definition of the Hamiltonian.
The obstruction persists because the mode identity carries a hidden
linear code accessible coherently to one quantum query, whereas each
classical coordinate supplies only one linear equation.

There need not be a singularity of the limiting pressure at any finite
positive inverse temperature. The same proof gives
\[
 Z_n(\beta)=2n\,[2\cosh(\beta L)]^n\mathbb E_qR,
 \qquad
 \lim_{n\to\infty}\frac{\log Z_n(\beta)}n
       =\log(2\cosh(\beta L))\quad(\beta>0).
\]
Indeed $\mathbb E_qR=1+O(e^{-cn})$. The displayed limiting pressure is
analytic throughout $\beta>0$, although full AA fails everywhere in
that interval. At $\beta=0$ the pressure equals $\log2$ exactly; no
claim about analyticity across negative inverse temperatures is needed.
